%=====================================================
\section{Tier 3.1: J-integral / dissipated-energy conservation}
\label{sec:tier3_1}

This section documents the governing equations and verification logic for
\texttt{tier3\_1\_jintegral.i} and the analytical reference implemented in
\texttt{verify\_tier3\_1.py}.

%=====================================================
\subsection{Test purpose}

Tier 3.1 verifies an \emph{integral} property of the cohesive law that the
pointwise tests (Tier 1.1--1.5) cannot reach by construction:
the energy dissipated per unit interface area at full decohesion equals
the input fracture toughness $G_{Ic}$.

In the small-strain limit this reduces to a one-line statement of Rice's
contour-independence theorem applied to the cohesive interface
$\Gamma_c$:
\begin{equation}
J \;=\;
\int_{0}^{\delta_n^{f}} t_n(\delta_n)\,d\delta_n
\;=\; \tfrac{1}{2}\,N\,\delta_n^{f}
\;=\; G_{Ic}.
\label{eq:Jintegral}
\end{equation}

What the test catches that Tier 1 does not:
\begin{enumerate}
\item a wrong interface measure $|\Gamma_c|$ in the surface (or line)
      integral,
\item a finite-strain reference--current configuration mix-up that scales
      the area mapping incorrectly while leaving the pointwise curve
      unchanged,
\item a sign error in the traction direction during integration.
\end{enumerate}
The first item is exercised here through an explicit \texttt{AreaPostprocessor}.
The second and third are exercised by the \emph{template} the test
provides; they become non-trivial in the finite-strain analogue.

%=====================================================
\subsection{Domain and loading}

Geometry, mesh, bulk material, and boundary conditions are identical to
Tier 1.1: a single-element mode-I monotonic pull on a $1\times 2$ unit
mesh with a horizontal cohesive interface $\Gamma_c=\{(x,y):\,0\le x\le1,\,
y=1\}$ split by \texttt{BreakMeshByBlockGenerator}.
The bottom face is clamped, the top face is constrained in $x$ and pulled
in $y$ by the same piecewise-linear function $u_*(t)$ used in Tier 1.1,
ramped past $u_*(1)=1.5\,\delta_n^{f}=0.03$ to guarantee complete
decohesion.

%=====================================================
\subsection{Two estimators of the cumulative work}

The verifier computes the cumulative work done per unit interface area by
the cohesive traction along the simulated trajectory two ways.

\paragraph{Trapezoidal-rule estimator.}
Sampled at the simulation time steps,
\begin{equation}
W_{\text{trap}}(t_k) \;=\;
\sum_{j=1}^{k} \tfrac{1}{2}\bigl(t_n^{(j)}+t_n^{(j-1)}\bigr)
                       \bigl(\delta_n^{(j)}-\delta_n^{(j-1)}\bigr).
\label{eq:Wtrap}
\end{equation}
This is the discrete, observation-only integral. It has trapezoidal-rule
error $O(\Delta t^{\,2})$ on smooth segments and $O(\Delta t)$ at the
bilinear-law kink at $\delta_n^{0}$.

\paragraph{Closed-form estimator.}
Using the running internal variable $\delta_n^{\max}$, the energy actually
dissipated by the bilinear law is
\begin{equation}
W_{\text{exact}}(\delta_n^{\max})
\;=\;
W_{\text{env}}(\delta_n^{\max})
\;-\; \tfrac{1}{2}\,(1-d)\,K\,\bigl(\delta_n^{\max}\bigr)^{2},
\label{eq:Wexact}
\end{equation}
where $W_{\text{env}}$ is the area swept under the bilinear envelope from
$0$ up to the running max, $d=d(\delta_n^{\max})$ is the damage variable,
and the second term subtracts the elastic energy still stored in the
interface at the current state. For the bilinear law this evaluates to
\begin{equation}
W_{\text{exact}}(\delta_n^{\max}) =
\begin{cases}
0,
& \delta_n^{\max} \le \delta_n^{0}, \\[6pt]
\dfrac{1}{2}\,K\,\delta_n^{0}\,
\dfrac{\delta_n^{\max}-\delta_n^{0}}{\delta_n^{f}-\delta_n^{0}}\,\delta_n^{f},
& \delta_n^{0}<\delta_n^{\max}<\delta_n^{f}, \\[12pt]
G_{Ic},
& \delta_n^{\max} \ge \delta_n^{f}.
\end{cases}
\label{eq:Wexact_piecewise}
\end{equation}
The third branch is exact: the bilinear law dissipates exactly $G_{Ic}$ at
full failure, so the closed-form estimator hits the target by construction.

\paragraph{Total work.}
Both per-area estimators are multiplied by the interface measure
$|\Gamma_c|$, reported by the in-MOOSE \texttt{AreaPostprocessor}:
\begin{equation}
W^{\,\text{tot}}_{\!\text{trap}}(t_k) = W_{\text{trap}}(t_k)\cdot|\Gamma_c|,
\qquad
W^{\,\text{tot}}_{\!\text{exact}}(t_k) = W_{\text{exact}}(\delta_n^{\max,(k)})\cdot|\Gamma_c|.
\end{equation}

%=====================================================
\subsection{Verification criterion}

The test passes when \emph{both} estimators land on $G_{Ic}\cdot|\Gamma_c|$
at the end of the simulation:
\begin{equation}
\bigl|\,W^{\,\text{tot}}_{\!\text{exact}}(t_{\text{end}}) - G_{Ic}\,|\Gamma_c|\,\bigr|
\;<\;10^{-9},
\qquad
\bigl|\,W^{\,\text{tot}}_{\!\text{trap}}(t_{\text{end}})  - G_{Ic}\,|\Gamma_c|\,\bigr|
\;<\;10^{-3}.
\end{equation}
The first inequality is a check on the cohesive law's damage-state
bookkeeping (it must arrive at $d=1$ with the correct
$\delta_n^{\max}\ge\delta_n^{f}$). The second is an integration cross-check
that catches any post-processing bug that might inflate or deflate
$\int t_n\,d\delta_n$ relative to the law itself.

%=====================================================
\subsection{Results}

For the parameters of Tier 1.1 ($K=5\times10^{3}$, $N=50$, $G_{Ic}=0.5$,
$|\Gamma_c|=1$) the verifier reports
\begin{align}
W^{\,\text{tot}}_{\!\text{exact}}(t_{\text{end}}) &= 0.500\,000\,000\,0
\quad\Rightarrow\quad
|\,\cdot - G_{Ic}\,| = 0,\\[2pt]
W^{\,\text{tot}}_{\!\text{trap}}(t_{\text{end}})  &= 0.499\,999\,952\,7
\quad\Rightarrow\quad
|\,\cdot - G_{Ic}\,| = 4.7\times 10^{-8}.
\end{align}
Both inside the respective tolerances by many orders of magnitude. The
trapezoidal residual is the kink-region quadrature error and shrinks as
$O(\Delta t)$ if the time step is refined.

%=====================================================
\subsection{Why the two estimators disagree \emph{during} the simulation}

The two estimators measure different quantities until full failure is
reached. The trap-rule estimator integrates the \emph{total} work done
by the traction, which equals
\[
W_{\text{trap}}(t) \;=\; W_{\text{diss}}(t) + U_{\text{elastic}}(t),
\]
i.e.\ dissipated energy plus the still-stored elastic energy in the
cohesive ``spring''. During the elastic loading branch
$\delta_n\le\delta_n^{0}$, the trap-rule estimator rises while the
closed-form estimator is identically zero (no damage yet, all stored).
Past the softening branch the trap-rule keeps rising and the closed-form
catches up; both meet at $G_{Ic}\,|\Gamma_c|$ when
$\delta_n\ge\delta_n^{f}$ because at that point $d=1$ and
$U_{\text{elastic}}=0$. The plot
\texttt{tier3\_1\_jintegral\_compare.png} shows this behaviour
explicitly.

\end{document}
